
Large language models have transformed code generation, achieving strong performance on benchmarks such as HumanEval and MBPP. Iterative repair loops---where models refine their output based on feedback---have emerged as a key technique for improving functional correctness. The dominant paradigm, exemplified by Self-Debug \citep{chen2023selfdebug}, relies on execution feedback: test results and stack traces guide the model toward working solutions.

Static analysis offers an alternative signal. Tools such as pylint and mypy detect potential issues---undefined variables, type mismatches, unreachable code---before execution. Unlike execution traces that show where code fails, static warnings explain why the code is problematic. This mechanistic feedback could enable more targeted repairs.

\textbf{The research gap.} Prior work has shown that static analysis improves code quality metrics. Blyth et al. (2025) demonstrated reductions in security issues (40\% to 13\%) and readability problems (80\% to 11\%) when integrating static feedback into LLM pipelines. However, no published study has measured the impact on functional correctness---the pass@k metrics that matter for code generation benchmarks.

\textbf{Research question.} Does integrating static analyzer feedback into iterative LLM code repair improve pass@k compared to execution-only feedback?

\textbf{This work.} This paper presents a methodology study that establishes groundwork for answering this question. A hypothesis validation framework with gate-based testing is developed, a pylint analysis pipeline for HumanEval is implemented, and baseline measurements are provided. The existence gate---testing whether static analyzers produce actionable warnings on benchmark code---reveals a methodological limitation: canonical solutions exhibit only 9.15\% warning rates, making them unsuitable proxies for LLM-generated code.

\textbf{Contributions:}

\begin{enumerate}
\item First baseline measurement of static warning rates on HumanEval canonical solutions (9.15\%, 23 warnings across 9 categories)
\item Validated pylint analysis pipeline for code generation research with reusable components
\item Gate-based hypothesis validation methodology that catches methodology limitations before full experiments
\item Identification of proxy limitation: canonical solutions do not represent LLM output for static analysis evaluation
\end{enumerate}

This work is framed as a methodology contribution. The hypothesis that static analysis improves pass@k remains untested, but the validated pipeline and baseline are provided for future work with LLM API access.
